\section{Contoh Lengkap RSA: Key Generation dan Enkripsi}

\subsection{Setup: Memilih Parameter}

Untuk memahami RSA secara konkret, kita akan menggunakan bilangan prima kecil. \textbf{Catatan}: Dalam praktik, gunakan prima 2048-bit atau lebih!

\textbf{Langkah 1: Pilih dua bilangan prima}
\begin{itemize}
  \item $p = 61$
  \item $q = 53$
\end{itemize}

\textbf{Verifikasi primalitas}: Kedua bilangan adalah prima.

\subsection{Key Generation Step-by-Step}

\textbf{Langkah 2: Hitung modulus $n$}
\begin{align*}
n &= p \times q \\
  &= 61 \times 53 \\
  &= 3233
\end{align*}

\textbf{Langkah 3: Hitung fungsi Euler $\phi(n)$}
\begin{align*}
\phi(n) &= (p-1)(q-1) \\
        &= 60 \times 52 \\
        &= 3120
\end{align*}

\textbf{Langkah 4: Pilih eksponen publik $e$}

Syarat: $1 < e < \phi(n)$ dan $\gcd(e, \phi(n)) = 1$

Pilih $e = 17$ (nilai umum: 3, 17, 65537)

Verifikasi: $\gcd(17, 3120) = 1$ $\checkmark$

\textbf{Langkah 5: Hitung eksponen privat $d$}

Cari $d$ sehingga $e \cdot d \equiv 1 \pmod{\phi(n)}$

Menggunakan Extended Euclidean Algorithm:
\begin{align*}
17 \cdot d &\equiv 1 \pmod{3120} \\
d &= e^{-1} \text{ mod } \phi(n) \\
  &= 2753
\end{align*}

\textbf{Verifikasi}:
\begin{align*}
17 \times 2753 &= 46801 \\
46801 \text{ mod } 3120 &= 1 \quad \checkmark
\end{align*}

\subsection{Kunci RSA yang Dihasilkan}

\begin{table}[htbp]
  \centering
  \begin{tabular}{ll}
    \toprule
    \textbf{Parameter} & \textbf{Nilai} \\
    \midrule
    Public Key $(e, n)$ & $(17, 3233)$ \\
    Private Key $(d, n)$ & $(2753, 3233)$ \\
    \midrule
    Prima $p$ & 61 (rahasia) \\
    Prima $q$ & 53 (rahasia) \\
    $\phi(n)$ & 3120 (rahasia) \\
    \bottomrule
  \end{tabular}
  \caption{Kunci RSA yang dihasilkan dari contoh.}
\end{table}

\subsection{Enkripsi Pesan}

\textbf{Plaintext}: $m = 123$ (harus $m < n$)

\textbf{Enkripsi}: $c = m^e \bmod n$

\begin{align*}
c &= 123^{17} \bmod 3233
\end{align*}

\textbf{Perhitungan menggunakan square-and-multiply}:

Eksponen $17 = 10001_2$ (binary)

\begin{table}[htbp]
  \centering
  \small
  \begin{tabular}{cccl}
    \toprule
    \textbf{Bit} & \textbf{Square} & \textbf{Multiply} & \textbf{Result} \\
    \midrule
    1 & - & $123^1$ & 123 \\
    0 & $123^2 = 15129 \bmod 3233$ & - & 855 \\
    0 & $855^2 = 731025 \bmod 3233$ & - & 2790 \\
    0 & $2790^2 = 7784100 \bmod 3233$ & - & 1084 \\
    1 & $1084^2 = 1175056 \bmod 3233$ & $\times 123$ & 2516 \\
    \bottomrule
  \end{tabular}
  \caption{Square-and-multiply untuk $123^{17} \bmod 3233$.}
\end{table}

\textbf{Ciphertext}: $c = 855$

\subsection{Dekripsi Pesan}

\textbf{Ciphertext}: $c = 855$

\textbf{Dekripsi}: $m = c^d \bmod n$

\begin{align*}
m &= 855^{2753} \bmod 3233
\end{align*}

Menggunakan square-and-multiply (eksponen besar, hanya hasil akhir):

\textbf{Plaintext yang didekripsi}: $m = 123$ $\checkmark$

\subsection{Verifikasi Kebenaran}

\begin{align*}
m &= (m^e)^d \bmod n \\
  &= m^{ed} \bmod n \\
  &= m^{ed \bmod \phi(n)} \bmod n \quad \text{(Euler's theorem)} \\
  &= m^1 \bmod n \\
  &= m
\end{align*}

Karena $ed = 17 \times 2753 = 46801 \equiv 1 \pmod{3120}$

\subsection{Contoh dengan Pesan Teks}

Untuk mengenkripsi teks, konversi ke angka terlebih dahulu:

\textbf{Plaintext}: "HI"
\begin{itemize}
  \item H = 72 (ASCII)
  \item I = 73 (ASCII)
\end{itemize}

\textbf{Enkripsi}:
\begin{align*}
c_1 &= 72^{17} \bmod 3233 = 1234 \\
c_2 &= 73^{17} \bmod 3233 = 2819
\end{align*}

\textbf{Ciphertext}: [1234, 2819]

\textbf{Dekripsi}:
\begin{align*}
m_1 &= 1234^{2753} \bmod 3233 = 72 \rightarrow \text{H} \\
m_2 &= 2819^{2753} \bmod 3233 = 73 \rightarrow \text{I}
\end{align*}

\textbf{Plaintext}: "HI" $\checkmark$

\subsection{Implementasi C untuk Contoh Ini}

\begin{lstlisting}[style=cstyle, caption={RSA dengan Parameter Kecil}]
#include <stdio.h>

// Modular exponentiation: (base^exp) % mod
long long mod_exp(long long base, long long exp, long long mod) {
    long long result = 1;
    base = base % mod;
    
    while (exp > 0) {
        if (exp % 2 == 1) {
            result = (result * base) % mod;
        }
        exp = exp >> 1;
        base = (base * base) % mod;
    }
    
    return result;
}

int main() {
    long long n = 3233;
    long long e = 17;
    long long d = 2753;
    
    // Plaintext
    long long m = 123;
    printf("Plaintext: %lld\n", m);
    
    // Encryption
    long long c = mod_exp(m, e, n);
    printf("Ciphertext: %lld\n", c);
    
    // Decryption
    long long decrypted = mod_exp(c, d, n);
    printf("Decrypted: %lld\n", decrypted);
    
    return 0;
}
\end{lstlisting}

\textbf{Output}:
\begin{verbatim}
Plaintext: 123
Ciphertext: 855
Decrypted: 123
\end{verbatim}

\begin{catatan}
Contoh ini menggunakan bilangan prima kecil untuk tujuan edukatif. Dalam praktik:
\begin{itemize}
  \item Gunakan prima 2048-bit atau lebih (RSA-2048, RSA-4096)
  \item Gunakan padding scheme seperti OAEP untuk keamanan
  \item Jangan enkripsi data panjang langsung dengan RSA (gunakan hybrid encryption)
  \item Gunakan library kriptografi yang sudah diaudit (OpenSSL, libsodium)
\end{itemize}
\end{catatan}

\subsection{Mengapa RSA Aman?}

Keamanan RSA bergantung pada kesulitan \textbf{faktorisasi bilangan besar}:

\begin{itemize}
  \item Publik tahu: $n = 3233$, $e = 17$
  \item Untuk menemukan $d$, penyerang harus:
    \begin{enumerate}
      \item Faktorkan $n = 3233$ menjadi $p \times q$
      \item Hitung $\phi(n) = (p-1)(q-1)$
      \item Hitung $d = e^{-1} \bmod \phi(n)$
    \end{enumerate}
  \item Untuk $n$ kecil (3233), ini mudah: $3233 = 61 \times 53$
  \item Untuk $n$ besar (2048-bit), ini \textbf{sangat sulit} dengan teknologi saat ini
\end{itemize}

\textbf{Contoh kompleksitas}:
\begin{itemize}
  \item RSA-768 (232 digit): difaktorkan tahun 2009 setelah 2 tahun komputasi
  \item RSA-2048 (617 digit): diperkirakan butuh jutaan tahun dengan teknologi saat ini
\end{itemize}
